\chapter{Физическая шкала двухузлового компакттона}
\label{ch:v6-spectral-transition-discrete-compacton-physical-scale-map-gate}

\section{Вопрос после локальной устойчивости}

Минимальная ветвь $\kappa=2\pi$ существует точно и не имеет расширяющей
флоке-моды в проверенных объёмах. Теперь требуется различить два утверждения:
\begin{enumerate}
 \item безразмерная орбита фиксирует фазу и число занятых узлов;
 \item теория предсказывает физическую длину, время и массу.
\end{enumerate}
Второе не следует автоматически из первого.

Строгий масштабный реестр проекта уже содержит три предупреждения. Гейт
\path{version4_absolute_scale_eft_validity_gate} показал, что формальное
планковское сопоставление нарушает режим собственного локального
heat-kernel разложения. Гейт
\path{version5_projector_superconnection_common_scale_gate} сохранил
независимые $f_2,f_0,\Lambda,\ell$. Наконец, гейт
\path{version6_single_thread_scale_hierarchy_branch_decision_gate} не
нашёл скрытой абсолютной длины или энергии.

Первичная литература согласуется с такой постановкой: решёточный и
временной шаги входят в непрерывный предел квантового автомата, а
квазиэнергия дискретной эволюции определена по модулю $2\pi$;
см. \path{arXiv:1212.2839} и \path{arXiv:1912.09768}. Эти факты не
запрещают фундаментальную решётку, но требуют отдельного закона для её
физической нормировки.

\section{Ранг размерных условий}

Обозначим пространственный шаг через $a$, временной шаг через $\Delta t$ и
энергию минимальной зоны через $E$. Если предельная скорость отождествлена
с $c$, то встречный сдвиг даёт
\begin{equation}
 \frac{a}{\Delta t}=c.
 \label{eq:v6-compacton-causal-scale}
\end{equation}
Собственная фаза $\pm i$ даёт минимальный модуль квазиэнергетического угла
$\vartheta=\pi/2$ и условное соотношение
\begin{equation}
 E\Delta t=\frac{\pi\hbar}{2}.
 \label{eq:v6-compacton-quasienergy-scale}
\end{equation}

Для переменных $(\log a,\log\Delta t,\log E)$ две формулы имеют матрицу
ранга два. Остаётся одномерное ядро
\begin{equation}
 (a,\Delta t,E)\longmapsto
 (\lambda a,\lambda\Delta t,E/\lambda).
 \label{eq:v6-compacton-scale-null-direction}
\end{equation}
Следовательно, скорость и собственная фаза не определяют абсолютную
единицу. Они фиксируют только обратную связь размера и энергии.

Для двухузлового диаметра $L=2a$ из
\eqref{eq:v6-compacton-causal-scale}--
\eqref{eq:v6-compacton-quasienergy-scale} получается
\begin{equation}
 \boxed{EL=\pi\hbar c},
 \qquad
 \frac{L}{\lambda_C}=\pi,
 \label{eq:v6-compacton-conditional-mass-size-product}
\end{equation}
где $\lambda_C=\hbar c/E$. Это новое условное безразмерное следствие, но
не абсолютная масса: любое $\lambda>0$ сохраняет
\eqref{eq:v6-compacton-conditional-mass-size-product}.

\section{Почему масштаб объединения не закрывает ядро}

Даже если принять проектное значение
\begin{equation}
 \Lambda_{\rm S2T}=\sqrt\pi\,10^{16}\ {\rm GeV},
\end{equation}
нужна дополнительная формула
\begin{equation}
 a=\alpha\frac{\hbar c}{\Lambda_{\rm S2T}}.
 \label{eq:v6-compacton-alpha-assignment}
\end{equation}
Коэффициент $\alpha$ текущим родителем не фиксирован. Три одинаково
совместимых примера дают:
\begin{center}
\begin{tabular}{@{}rrrr@{}}
\toprule
$\alpha$ & $a$, м & $L=2a$, м & $E$, ГэВ \\
\midrule
$0.1$ & $1.1133\cdot10^{-33}$ & $2.2266\cdot10^{-33}$ & $2.7842\cdot10^{17}$ \\
$1$ & $1.1133\cdot10^{-32}$ & $2.2266\cdot10^{-32}$ & $2.7842\cdot10^{16}$ \\
$\pi$ & $3.4975\cdot10^{-32}$ & $6.9951\cdot10^{-32}$ & $8.8623\cdot10^{15}$ \\
\bottomrule
\end{tabular}
\end{center}
Во всех строках $EL/(\hbar c)=\pi$. Совпадение произведения не выбирает
строку и потому не превращает $\Lambda_{\rm S2T}$ в шаг автомата.

\section{Модулярность квазиэнергии}

Для дискретной эволюции углы
\begin{equation}
 \vartheta_n=\pm\frac\pi2+2\pi n,
 \qquad n\in\mathbb Z,
\end{equation}
задают ту же собственную фазу. Выбор главной зоны оставляет пару
$\pm\pi/2$, но физическая энергия всё равно требует $\Delta t$ и
опорного вакуумного уровня. Поэтому фаза $\pm i$ является точной
квазиэнергетической меткой орбиты, но не абсолютной массой сама по себе.

\section{Несовместимость с фиксированным непрерывным пределом}

В исходном нелинейном дискретном родителе контролируемый дираковский предел
строился масштабированием
\begin{equation}
 \kappa=ga,
 \qquad a\longrightarrow0,
\end{equation}
при конечной физической связи $g$. Компактон, напротив, требует
\begin{equation}
 \kappa=2\pi.
\end{equation}
Чтобы сохранить эту ветвь при $a\to0$, пришлось бы положить
\begin{equation}
 g=\frac{2\pi}{a}\longrightarrow\infty.
 \label{eq:v6-compacton-divergent-continuum-coupling}
\end{equation}
Следовательно, найденный объект является собственно решёточным
двухузловым состоянием, а не автоматически контролируемым солитоном
фиксированной континуальной теории. При $c=1$ сохраняется ещё одно
условное отношение $E/g=1/4$, но обе размерные величины масштабируются
совместно как $1/a$.

\begin{theorem}[Масштабный статус компакттона]
Точная опора на двух узлах и фаза $\pm i$ фиксируют условное произведение
$EL=\pi\hbar c$, но оставляют одно непрерывное масштабное направление.
Ни $\Lambda_{\rm S2T}$, ни прежняя планковская нормировка, ни общий
оператор суперсвязности не дают выведенного уравнения для $a$. Поэтому
абсолютные размер и масса компакттона не предсказаны.
\end{theorem}

\begin{nogocriterion}
Нельзя подставлять $a=\hbar c/\Lambda_{\rm S2T}$ или $a=\ell_P$ без
отдельного родительского закона. Нельзя также переносить локальную
флоке-устойчивость непосредственно в континуальную теорию: сохранение
$\kappa=2\pi$ требует расходящейся связи $g\sim1/a$.
\end{nogocriterion}

\section{Следующий гейт}

Следующий гейт
\path{version6_spectral_transition_discrete_compacton_dynamical_capture_gate}
должен проверить, возникает ли компактон из открытого множества обычных
локализованных начальных состояний либо существует только как специально
приготовленная орбита.

Нулевая гипотеза: нормосохраняющая дискретная динамика не захватывает
общий пакет в двухузловую орбиту, а лишь рассеивает его или оставляет
квазипериодическим. Стоп-критерий: отсутствие воспроизводимой области
захвата закрывает компактон как механизм рождения, сохраняя его только как
точное решёточное состояние.

\begin{protocol}
\begin{enumerate}
 \item точная квазиэнергетическая фаза: \textbf{$\pm\pi/2$ по модулю $2\pi$};
 \item двухузловая опора: \textbf{да};
 \item ранг размерных условий: \textbf{два};
 \item число свободных масштабных направлений: \textbf{одно};
 \item условное произведение: \textbf{$EL=\pi\hbar c$};
 \item связь $a$ с $\Lambda_{\rm S2T}$: \textbf{не выведена};
 \item фиксированный-$g$ непрерывный предел компакттона: \textbf{не проходит};
 \item абсолютный размер: \textbf{не предсказан};
 \item абсолютная масса: \textbf{не предсказана};
 \item R5: \textbf{не закрыт}.
\end{enumerate}
\end{protocol}

Машинный сертификат сохранён в
\path{s2t_v6_spectral_transition_discrete_compacton_physical_scale_map_gate_results.json}.